\documentclass[12pt,a4paper]{article}

\usepackage[T1]{fontenc}
\usepackage[utf8]{inputenc}
\usepackage[brazil]{babel}
\usepackage{lmodern}
\usepackage{geometry}
\usepackage{enumitem}
\usepackage{amssymb}
\usepackage[hidelinks]{hyperref}

\geometry{margin=2.5cm}
\setlist[itemize]{leftmargin=1.5em}
\setlist[enumerate]{leftmargin=1.8em,itemsep=0.5em}

\title{Checklist para execução da metodologia}
\author{Luiz Gabriel Bezerra da Conceição}
\date{Outubro de 2026}

\begin{document}

\maketitle

\section*{Uso deste checklist}

Este checklist organiza o que preciso concluir antes, durante e depois da intervenção pedagógica. A aplicação envolverá uma sequência didática de oito encontros, com uso da plataforma interativa, mediação docente e produções individuais dos estudantes. As conclusões deverão se limitar às evidências produzidas na turma participante.

\newcommand{\checkitem}[1]{\item[$\square$] #1}

\section*{Passo 1. Confirmar o campo de aplicação}

\begin{itemize}
    \checkitem{Definir a escola em que a intervenção será realizada.}
    \checkitem{Definir a turma participante.}
    \checkitem{Definir a série e o componente curricular.}
    \checkitem{Confirmar o número aproximado de estudantes.}
    \checkitem{Confirmar o calendário e reservar os oito encontros de aproximadamente 50 minutos.}
    \checkitem{Verificar a infraestrutura, especialmente cerca de 20 computadores do laboratório, Chromebooks disponíveis, internet, projeção e plano alternativo para falhas técnicas.}
    \checkitem{Confirmar que cada estudante terá acesso individual à plataforma, preservando os registros individuais previstos para a pesquisa.}
\end{itemize}

\section*{Passo 2. Regularizar os procedimentos éticos e institucionais}

\begin{itemize}
    \checkitem{Confirmar, com a instituição e o orientador, quais autorizações e apreciações éticas são exigidas.}
    \checkitem{Preparar o projeto, os instrumentos e os documentos necessários para submissão ao comitê de ética aplicável, quando exigido.}
    \checkitem{Obter autorização da escola ou da rede de ensino.}
    \checkitem{Preparar os termos de consentimento dos responsáveis e de assentimento dos estudantes, quando aplicáveis.}
    \checkitem{Informar que a participação é voluntária e não interfere em nota, frequência ou acesso às aulas.}
    \checkitem{Planejar atividade pedagógica equivalente para estudantes que não participarem da pesquisa.}
    \checkitem{Definir onde os dados serão guardados, quem terá acesso, por quanto tempo serão mantidos e como serão descartados.}
\end{itemize}

\section*{Passo 3. Preparar a sequência, a plataforma e os instrumentos}

\begin{itemize}
    \checkitem{Revisar os cenários, as explicações e os controles da plataforma de acordo com a turma.}
    \checkitem{Revisar a sequência dos oito encontros.}
    \checkitem{Associar cada atividade a um objetivo de aprendizagem.}
    \checkitem{Preparar a matriz de comparabilidade entre diagnóstico e instrumento final.}
    \checkitem{Submeter os instrumentos à leitura do orientador e, quando possível, de um professor com experiência no Ensino Médio.}
    \checkitem{Testar a clareza dos enunciados com pessoa de perfil próximo ao público, quando isso for viável e permitido.}
    \checkitem{Preparar códigos para anonimizar as produções dos participantes.}
\end{itemize}

\section*{Passo 4. Aplicar o diagnóstico inicial}

\begin{itemize}
    \checkitem{Aplicar o diagnóstico individualmente, antes do primeiro uso orientado da plataforma.}
    \checkitem{Explicar que o instrumento não vale nota e que respostas em branco podem ser entregues.}
    \checkitem{Não permitir consulta à plataforma, à internet, ao material de aula ou aos colegas durante o diagnóstico.}
    \checkitem{Usar apenas o código do participante no instrumento.}
    \checkitem{Registrar tempo de aplicação, dúvidas sobre o enunciado, interrupções e instrumentos incompletos.}
    \checkitem{Utilizar o instrumento: \href{../../03_pesquisa/instrumentos/diagnostico_inicial_v01/instrumento_diagnostico_inicial_v01.pdf}{\textbf{Diagnóstico inicial, versão 01 (PDF)}}.}
\end{itemize}

\section*{Passo 5. Executar a sequência didática}

\begin{itemize}
    \checkitem{Realizar os oito encontros previstos: dados e tipos de tarefa; regressão linear; regressão logística; \textit{k}-vizinhos mais próximos; árvore de decisão; \textit{k}-means; comparação entre modelos; síntese final.}
    \checkitem{Em cada atividade, solicitar que os estudantes prevejam, experimentem, observem e expliquem.}
    \checkitem{Fazer registros individuais antes e depois da exploração nos encontros 2, 4 e 6.}
    \checkitem{Ao final de cada encontro, realizar uma síntese coletiva sem fornecer uma resposta para ser apenas copiada.}
    \checkitem{Registrar tempo utilizado, dificuldades recorrentes, intervenções realizadas, adaptações e condições de infraestrutura.}
    \checkitem{Não usar como dados da pesquisa as respostas digitadas no site, pois elas têm apenas função formativa e não ficam armazenadas.}
\end{itemize}

\section*{Passo 6. Aplicar o instrumento final e o questionário de percepção}

\begin{itemize}
    \checkitem{Aplicar individualmente o núcleo final com questões comparáveis ao diagnóstico.}
    \checkitem{Aplicar o desafio de transferência com uma situação diferente das usadas nas aulas.}
    \checkitem{Aplicar o questionário breve de percepção sobre clareza, dificuldade e utilidade das atividades e da plataforma.}
    \checkitem{Manter a análise do questionário de percepção separada da análise de aprendizagem.}
    \checkitem{Registrar ausências, respostas em branco e instrumentos incompletos.}
\end{itemize}

\section*{Passo 7. Organizar e analisar os dados}

\begin{itemize}
    \checkitem{Substituir nomes por códigos e guardar separadamente qualquer dado identificável.}
    \checkitem{Ler integralmente as respostas antes de classificar o corpus.}
    \checkitem{Aplicar a rubrica de níveis 0, 1 e 2 nas dimensões: dados e finalidade; alteração e consequência; representação e conceito; possibilidades e limites.}
    \checkitem{Registrar uma justificativa curta para cada classificação.}
    \checkitem{Comparar apenas as questões planejadas como equivalentes entre diagnóstico e instrumento final.}
    \checkitem{Produzir tabelas de transição por dimensão e frequências com os respectivos denominadores.}
    \checkitem{Analisar o desafio de transferência separadamente.}
    \checkitem{Usar os registros docentes para contextualizar as respostas, sem completar ou alterar a resposta individual do estudante.}
    \checkitem{Solicitar, quando viável, que um segundo avaliador classifique parte das produções e registrar os critérios discutidos.}
\end{itemize}

\section*{Passo 8. Redigir e interpretar os resultados}

\begin{itemize}
    \checkitem{Apresentar mudanças, permanências e dificuldades por dimensão de análise.}
    \checkitem{Incluir excertos anonimizados que exemplifiquem diferentes formas de explicação.}
    \checkitem{Apresentar o desafio de transferência e a percepção dos estudantes em seções próprias.}
    \checkitem{Descrever as condições reais da aplicação e as adaptações realizadas.}
    \checkitem{Não atribuir as mudanças exclusivamente à plataforma.}
    \checkitem{Não generalizar os resultados para todos os estudantes do Ensino Médio.}
\end{itemize}

\end{document}
